\documentclass[12pt]{article}
\usepackage{graphicx}
\usepackage{amsmath}
\usepackage{amsfonts}
\usepackage{multirow}
\usepackage[margin=0.6in]{geometry}  

\begin{document}
The self-attention mechanism can be written as:
\begin{align*} 
\operatorname{Self}&\operatorname{Attention}(X; W_q, W_k, W_v) \\
&= \operatorname{softmax}\!\left(\frac{(XW_q)(XW_k)^\top}{\sqrt{d_k}}\right)(XW_v)
\end{align*}

The relational-attention mechanism can be written as:
\begin{align*} 
\operatorname{Relational}&\operatorname{Attention}(X; W_q, W_k, V) \\
&= \operatorname{softmax}\!\left(\frac{(XW_q)(XW_k)^\top}{\sqrt{d_k}}\right)V
\end{align*}

The exclusive or (XOR) task can be described by:
\begin{table}[h]
\centering
\begin{tabular}{c c | c}
\hline
$A$ & $B$ & $A \oplus B$ \\
\hline
0 & 0 & 0 \\
0 & 1 & 1 \\
1 & 0 & 1 \\
1 & 1 & 0 \\
\hline
\end{tabular}
\caption{Truth table for XOR.}
\end{table}

The pairwise transitive inference task can be described by:
\begin{table}[h]
\centering
\begin{tabular}{c c c | c}
\hline
Obs.? & $x_i$ & $x_j$ & $x_i > x_j$ \\
\hline
\multirow{4}{*}{Yes} 
  & A & B & 1 \\
  & B & C & 1 \\
  & B & A & 0 \\
  & C & B & 0 \\
\hline
\multirow{2}{*}{No} 
  & A & C & 1 \\
  & C & A & 0 \\
\hline
\end{tabular}
\caption{Truth table for transitive inference.}
\end{table}

\end{document}